\chapter{Model zarządzania uprawnieniami i analiza wartości}
\label{ch:governance}

Rozdział ten łączy dwa wątki. Najpierw domyka ostatni punkt centralizacji
w~systemie UniVerify --- pojedynczy klucz właściciela kontraktu --- zastępując go
modelem zarządzania opartym na portfelu wielopodpisowym i~mechanizmie opóźnienia
czasowego. Następnie, mając już kompletny obraz kosztów (rozdz.~\ref{ch:l2})
i~własności bezpieczeństwa systemu, przeprowadza analizę wartości: porównuje
UniVerify z~podejściem tradycyjnym oraz z~infrastrukturą klucza publicznego i~stawia
pytanie, czy przewaga blockchaina nad dojrzalszymi alternatywami uzasadnia jego
zastosowanie.

\section{Problem centralizacji uprawnień}
\label{sec:gov-problem}

W~systemie podstawowym zbiorem autoryzowanych wystawców zarządza jedno konto
zewnętrzne (\emph{externally owned account}, EOA) pełniące rolę właściciela
kontraktu \code{DiplomaRegistry}. Właściciel może dodawać i~usuwać wystawców oraz
uczelnie. Taka konstrukcja jest prosta, lecz wprowadza pojedynczy punkt awarii
o~znaczeniu krytycznym dla całego modelu zaufania.

Przejęcie tego jednego klucza pozwala napastnikowi w~jednej transakcji i~ze
skutkiem natychmiastowym przepisać zbiór uprawnionych wystawców: autoryzować
własny adres i~wydawać fałszywe, lecz formalnie poprawne dyplomy, albo odebrać
uprawnienia legalnej uczelni. Co istotne, zmiana taka jest \emph{cicha} --- dalej
w~tym rozdziale termin ten oznacza zmianę nieodróżnialną od uprawnionej operacji
administracyjnej i~niewykrywalną, zanim nabierze mocy. Osłabia to zasadniczą tezę
pracy: rejestr, który
jeden podmiot może niepostrzeżenie zmodyfikować, nie jest w~praktyce
istotnie odporniejszy na manipulację niż baza danych pod kontrolą tego samego
podmiotu.

Problem nie dotyczy przy tym pojedynczego dyplomu --- te są zakotwiczone
w~niezmiennych korzeniach Merkle --- lecz \emph{warstwy uprawnień} decydującej
o~tym, czyje korzenie są uznawane za wiarygodne. To właśnie ta warstwa wymaga
wzmocnienia.

\section{Model zarządzania: multisig + timelock}
\label{sec:gov-model}

Proponowany model przenosi własność rejestru z~pojedynczego EOA na dwuwarstwową
konstrukcję, nie zmieniając przy tym ani jednego bajtu kontraktu
\code{DiplomaRegistry}. Instalacja odbywa się wyłącznie przez istniejący,
dwustopniowy mechanizm przekazania własności
(\code{transferOwnership} $\rightarrow$ \code{acceptOwnership}), dzięki czemu
pomiary kosztów z~rozdziału~\ref{ch:l2} pozostają w~mocy.

\paragraph{Warstwa opóźnienia.} Nowym właścicielem rejestru zostaje kontrakt
\code{TimelockController} z~biblioteki OpenZeppelin~\cite{oz-contracts}. Każda
operacja administracyjna musi najpierw zostać \emph{zaplanowana}, po czym może być
\emph{wykonana} dopiero po upływie zadanego opóźnienia (\code{minDelay}). Rola
wykonawcy jest otwarta --- po upływie opóźnienia zaplanowaną operację może wywołać
dowolny podmiot --- natomiast wyłącznym \emph{proposerem} i~\emph{cancellerem} jest
portfel wielopodpisowy opisany niżej. Adres administratora timelocka ustawiono na
\code{address(0)}, co trwale usuwa administracyjną furtkę.

\paragraph{Warstwa progu.} Jedynym podmiotem uprawnionym do planowania
i~anulowania operacji jest \code{RegistryMultisig} --- portfel $m$-z-$n$ napisany
na potrzeby tej pracy, w~konwencji klasycznego \code{MultiSigWallet} firmy
ConsenSys~\cite{consensys-multisig}. Wybór implementacji własnej zamiast
audytowanego standardu branżowego (Safe) jest podyktowany celem pomiarowym:
Safe zbiera podpisy \emph{poza łańcuchem} i~wykonuje operację jedną transakcją,
przez co narzut progu staje się w~pomiarze niewidoczny. Konstrukcja użyta tutaj
zapisuje na łańcuchu każdy krok --- zgłoszenie, potwierdzenia, wykonanie ---
i~dzięki temu pozwala zmierzyć pełny, jawny koszt zarządzania
(sekcja~\ref{sec:gov-koszt}). Cena tej przejrzystości jest jednak realna: koszt
zgłoszenia operacji jest o~rząd wielkości wyższy niż w~Safe, a~implementacja
własna nie przeszła audytu i~nosi ograniczenie rezydualne opisane niżej. Do
wdrożenia produkcyjnego właściwym wyborem jest Safe
(rozdz.~\ref{ch:podsumowanie}); liczby z~sekcji~\ref{sec:gov-koszt} należy wówczas
czytać jako \emph{górne ograniczenie} narzutu. Transakcja wymaga zebrania $m$~potwierdzeń
spośród $n$~bieżących właścicieli, zanim zostanie wykonana. Portfel jest
\emph{samorządny}: operacje \code{addOwner}, \code{removeOwner}
i~\code{changeThreshold} może wywołać wyłącznie sam kontrakt (poprzez wywołanie
\code{onlySelf}), a~więc tylko za zgodą progu. Umożliwia to rotację skompromitowanego
sygnatariusza bez wymiany całej infrastruktury.

\paragraph{Przyjęte parametry.} Model jest sparametryzowany trójką
$(\code{minDelay}, m, n)$ i~bez podania jej wartości nie da się ocenić żadnej
z~gwarancji omawianych niżej. Jako konfigurację docelową przyjęto
$\code{minDelay} = 48$~godzin przy progu $m=3$ z~$n=5$. Dobór opóźnienia jest
kompromisem: musi być na tyle długie, by niezależny obserwator zdążył zauważyć
zdarzenie \code{CallScheduled} i~zaalarmować sygnatariuszy w~cyklu pracy
instytucji (dwie doby obejmują weekend), a~zarazem na tyle krótkie, by rutynowa
zmiana listy wystawców nie paraliżowała administracji. Próg $3$ z~$5$ toleruje
niedostępność dwóch sygnatariuszy bez utraty żywotności i~wymaga przejęcia trzech
niezależnych kluczy do złamania bezpieczeństwa.

Pomiary kosztu w~sekcji~\ref{sec:gov-koszt} wykonano w~konfiguracji minimalnej
$m=n=2$; narzut rośnie z~$n$, ponieważ \code{execute} przelicza liczbę
potwierdzeń, iterując po bieżącym zbiorze właścicieli. We wdrożeniu testowym na
sieci Sepolia użyto $\code{minDelay} = 120$~sekund --- wartości wyłącznie
technicznej, pozwalającej przejść pełną ścieżkę zarządzania w~jednej sesji
testowej i~nieodpowiadającej żadnemu założeniu bezpieczeństwa.

\paragraph{Kto trzyma klucze.} Model nie \emph{usuwa} zaufania, lecz je
\emph{rozprasza} --- i~uczciwość wymaga wskazania, dokąd. Ktoś musi dysponować
$n$~kluczami, a~sposób wyboru i~wymiany tego zbioru pozostaje poza łańcuchem, więc
to on wyznacza faktyczny korzeń zaufania systemu. Naturalnym kandydatem jest
konsorcjum uczelni, w~którym każda instytucja wnosi jednego sygnatariusza:
zbiór jest wtedy jawny, a~żadna pojedyncza uczelnia ani organ nadzoru nie
kontroluje progu. Rozwiązaniem alternatywnym jest powierzenie zbioru organowi
centralnemu, co jednak odtwarza pojedynczy punkt zaufania w~nowym miejscu
i~znosi większość korzyści z~całej konstrukcji. Praca nie rozstrzyga tego
wyboru --- jest on instytucjonalny, nie techniczny --- lecz wskazuje, że bez
niego twierdzenie o~,,usunięciu pojedynczego punktu zaufania'' pozostaje
niepełne: przenosi je z~jednego klucza na regulamin zbioru sygnatariuszy.

Złożenie obu warstw daje następującą własność: żadna zmiana uprawnień nie jest
ani jednoosobowa, ani natychmiastowa. Wymaga $m$~podpisów \emph{oraz} przetrwania
okna opóźnienia, w~którym pozostaje publicznie widoczna i~odwracalna
(rys.~\ref{fig:governance}).

\begin{figure}[H]
\centering
\begin{tikzpicture}[node distance=12mm and 22mm,
                    blok/.append style={minimum width=35mm},
                    blokszary/.append style={minimum width=35mm},
                    blokakcent/.append style={minimum width=35mm}]
  \node[blokakcent] (ms) {\code{RegistryMultisig}\\próg $m$-z-$n$};
  \node[blok, above=of ms] (sig) {sygnatariusze ($n$ portfeli)};
  \node[blokakcent, right=of ms] (tl) {\code{TimelockController}\\opóźnienie \code{minDelay}};
  \node[blokszary, right=of tl] (reg) {\code{DiplomaRegistry}\\(owner = timelock)};
  \draw[strz] (sig) -- node[etyk, right=1mm] {\code{submit}, \code{confirm}$\times m$} (ms);
  \draw[strz] (ms) -- node[etyk, above=2mm] {\code{schedule}} (tl);
  \draw[strz] (tl) -- node[etyk, above=2mm] {\code{execute}\\{\footnotesize po opóźnieniu}} (reg);
  \draw[strz, dashed] (ms.south) to[out=-35, in=-145]
    node[etyk, below=2mm] {\code{cancel} w~oknie opóźnienia} (tl.south);
\end{tikzpicture}
\caption[Model zarządzania: multisig + timelock]{Model zarządzania: każda operacja administracyjna wymaga progu
$m$~podpisów, a~następnie przetrwania publicznego okna opóźnienia, w~którym może
zostać anulowana. Zaplanowanie operacji emituje zdarzenie \code{CallScheduled},
dzięki czemu zmiana jest publicznie widoczna, zanim nabierze mocy.}
\label{fig:governance}
\end{figure}

\section{Model zagrożeń}
\label{sec:gov-zagrozenia}

Tabela~\ref{tab:gov-threat} zestawia model pojedynczego EOA z~proponowanym modelem
$m$-z-$n$ z~opóźnieniem. Kluczowa różnica nie polega na tym, że nowego modelu
\emph{nie da się} zaatakować, lecz na tym, że atak przestaje być cichy,
natychmiastowy i~nieodwracalny.

\begin{table}[H]
\centering
\caption{Porównanie modeli zagrożeń warstwy uprawnień.}
\label{tab:gov-threat}
\begin{tabular}{@{}lcc@{}}
\toprule
Własność & Pojedynczy EOA & Multisig 3-z-5 + timelock 48~h \\
\midrule
Wyciek jednego klucza przepisuje uprawnienia & tak & nie --- wymaga 3~podpisów \\
Zmiana uprawnień cicha i~natychmiastowa & tak & nie --- planowana i~opóźniona \\
Zmiana wykrywalna, zanim zadziała & nie & tak --- w~oknie 48~h od \code{CallScheduled} \\
Złośliwa zmiana odwracalna & nie & tak --- \code{cancel} w~tym samym oknie \\
Odtworzenie po kompromitacji sygnatariusza & --- & tak --- \code{removeOwner}/\code{addOwner} \\
Żywotność przy niedostępnym sygnatariuszu & --- & toleruje 2~niedostępnych z~5 \\
\bottomrule
\end{tabular}
\end{table}

Trzy środkowe wiersze są sednem argumentu tej pracy. Zaplanowanie operacji emituje
na łańcuchu zdarzenie \code{CallScheduled}, a~opóźnienie \code{minDelay} tworzy okno,
w~którym każdy obserwator może wykryć nadchodzącą zmianę, zanim ta nabierze mocy.
W~tym samym oknie próg sygnatariuszy może operację anulować. Rejestr uprawnień staje
się więc \emph{publicznym, dopisywalnym i~odpornym na ciche wsteczne przepisanie}
logiem --- dokładnie tą własnością, do której wraca analiza wartości
w~sekcji~\ref{sec:gov-porownanie}.

\paragraph{Cena opóźnienia: reakcja awaryjna.} Opóźnienie chroni przed cichą
zmianą uprawnień, lecz w~jednym scenariuszu działa przeciw obrońcy. Gdy dojdzie
do przejęcia klucza \emph{wystawcy}, odebranie mu uprawnień wymaga zebrania
$m$~podpisów \emph{oraz} przeczekania \code{minDelay} --- przy przyjętych 48~h
napastnik dysponuje dwiema dobami na publikowanie kolejnych fałszywych korzeni.
Sytuację pogarsza brak unieważniania na poziomie partii: kontrakt udostępnia
wyłącznie \code{revokeFromBatch}, działające na pojedynczym dokumencie, więc
wycofanie sfałszowanej partii 10~tys.\ dyplomów wymagałoby 10~tys.\ transakcji.

Jest to realna luka projektowa, nie tylko brak opisu. Wskazać można dwie ścieżki
jej domknięcia, obie pozostawione poza zakresem implementacyjnym pracy. Pierwsza
to \code{revokeBatchRoot} --- operacja unieważniająca cały korzeń jednym zapisem,
tania i~prosta, lecz gruboziarnista: cofa również dyplomy wydane w~tej partii
uczciwie. Druga to wydzielenie roli \emph{strażnika} (\emph{guardian})
z~uprawnieniem wyłącznie do anulowania zaplanowanych operacji i~wstrzymania
wystawcy, bez prawa zmiany rejestru --- ścieżka szybka, bo pomijająca opóźnienie,
a~zarazem niegroźna, bo z~definicji nie może niczego dodać. Druga jest właściwsza:
zachowuje asymetrię, w~której działania \emph{rozszerzające} uprawnienia są powolne
i~jawne, a~\emph{ograniczające} --- natychmiastowe.

\paragraph{Ograniczenie rezydualne.} Operacja \code{removeOwner} nie czyści
potwierdzeń usuniętego właściciela na transakcjach nadal oczekujących. Jeśli ten sam
adres zostanie później ponownie dodany, jego stare potwierdzenie zaczyna liczyć się
na nowo. Jest to standardowy kompromis wzorca ConsenSys, unikający nieograniczonej
iteracji po zbiorze oczekujących transakcji przy każdym usunięciu. Operacyjnie
należy unikać ponownego dodawania byłego właściciela, gdy istnieją wrażliwe
transakcje oczekujące.

\section{Koszt zarządzania}
\label{sec:gov-koszt}

Narzut modelu zmierzono testem \code{GovernanceGasTest}
(\code{forge test --gas-report}) w~konfiguracji minimalnej $m=n=2$; dla progu
docelowego 3-z-5 narzut jest wyższy o~koszt iteracji po większym zbiorze
właścicieli w~\code{execute}. Wszystkie wdrożenia kontraktów oraz jednorazowe
przekazanie własności realizowane są w~\code{setUp()}, dzięki czemu ciało każdego
testu mierzy dokładnie jedną operację onboardu wystawcy, a~raportowana przez Foundry
suma na test (wyłączająca \code{setUp}) jest liczbą podaną bezpośrednio, niewymagającą
dekompozycji.

Linią bazową sprzed wprowadzenia zarządzania jest pojedyncze wywołanie
\code{onboardIssuer\allowbreak AndUniversity} na rejestrze będącym własnością zwykłego EOA:
84\,246 gazu. Ta sama operacja przeprowadzona pełną ścieżką zarządzaną ---
\code{multisig.submit} + \code{multisig.confirm} + \code{multisig.execute}
(planuje na timelocku) oraz \code{timelock.execute} (wywołuje rejestr po
opóźnieniu) --- kosztuje 824\,254 gazu, co daje narzut na zmianę uprawnień rzędu
$824\,254 - 84\,246 = 740\,008$ gazu.

Tabela~\ref{tab:gov-cost} pokazuje, gdzie trafia gaz ścieżki zarządzanej. Dominuje
\code{submit}, który zapisuje na łańcuchu pełne dane zaplanowanej operacji.

\begin{table}[H]
\centering
\caption{Rozbicie kosztu zarządzanego onboardu (wykonanie on-chain).}
\label{tab:gov-cost}
\begin{tabular}{@{}lrl@{}}
\toprule
Wywołanie & Gaz & Uwaga \\
\midrule
\code{RegistryMultisig.submit}   & 487\,647 & dominuje --- zapis danych operacji \\
\code{RegistryMultisig.confirm}  &  52\,213 & niezależne od danych \\
\code{RegistryMultisig.execute}  & 160\,119 & wywołuje \code{timelock.schedule} \\
\code{TimelockController.execute} & 99\,240 & wywołuje rejestr \\
\midrule
Suma & 799\,219 & samo wykonanie on-chain \\
\bottomrule
\end{tabular}
\end{table}

Suma czterech wywołań (799\,219 gazu) jest niższa od kosztu całego testu
(824\,254 gazu) o~narzut ciała testu --- przygotowanie i~kodowanie ABI
argumentów wywołań, około 3\% --- który nie jest kosztem on-chain. Porównanie
z~linią bazową używa spójnie sum na~test po~obu stronach, więc narzut ten
znosi się w~różnicy jedynie częściowo i~podana wartość 740\,008 gazu jest
oszacowaniem ostrożnym (z~góry).

Istotne jest przy tym, że zmiany uprawnień (dodanie lub usunięcie wystawcy) są
\emph{rzadkie}. Narzut 740\,008 gazu ponoszony jest raz na wystawcę i~amortyzuje się
na wszystkich dyplomach, które ten wystawca później wyda. Koszt rejestracji
pojedynczego dyplomu, zmierzony w~rozdziale~\ref{ch:l2}, pozostaje więc
niezmieniony.

\section{Porównanie podejść: status quo, PKI, blockchain}
\label{sec:gov-porownanie}

Mając kompletny obraz kosztów i~gwarancji, można postawić pytanie zasadnicze: co
właściwie daje blockchain wobec tańszych i~dobrze zrozumianych alternatyw?
Poniżej zestawiono trzy podejścia do weryfikacji dyplomu:

\begin{description}
  \item[Status quo] --- weryfikator kontaktuje się z~uczelnią (telefon, e-mail,
        formularz) i~prosi o~potwierdzenie. Dominujący dziś sposób.
  \item[PKI] --- uczelnia podpisuje cyfrowo dokument dyplomu kluczem
        certyfikowanym w~infrastrukturze klucza publicznego; weryfikator sprawdza
        podpis offline wobec certyfikatu wystawcy.
  \item[PKI + CT] --- to samo, uzupełnione o~publiczny, wyłącznie dopisywalny
        log wystawionych certyfikatów wystawców w~modelu \emph{Certificate
        Transparency}~\cite{rfc6962} (sekcja~\ref{sec:pki}).
  \item[Blockchain (UniVerify)] --- korzenie Merkle partii oraz zbiór uprawnionych
        wystawców publikowane są w~kontrakcie na publicznym blockchainie.
\end{description}

Tabela~\ref{tab:gov-value} zestawia je względem sześciu własności.

\begin{table}[H]
\centering
\caption{Macierz wartości: cztery podejścia wobec sześciu własności.}
\label{tab:gov-value}
\begin{tabular}{@{}lcccc@{}}
\toprule
Własność & \makecell{Status\\quo} & PKI & \makecell{PKI\\+ CT} & UniVerify \\
\midrule
Weryfikacja bez kontaktu z~wystawcą                & nie & tak & tak & tak \\
Niski koszt jednostkowej weryfikacji              & nie & tak & tak & tak \\
Działa, gdy wystawca przestał istnieć              & nie & częściowo & częściowo & tak \\
Publiczny, dopisywalny log uprawnień do~wydawania & nie & nie & tak & tak \\
Odporność na ciche wsteczne odebranie uprawnień   & --- & nie & tak & tak \\
Jeden widok logu bez warstwy dodatkowej           & --- & --- & nie & tak \\
\bottomrule
\end{tabular}
\end{table}

Ocena „częściowo'' w~trzecim wierszu wymaga komentarza. Podpis cyfrowy
pozostaje weryfikowalny offline tak długo, jak długo daje się zwalidować
łańcuch certyfikatów: po zniknięciu wystawcy certyfikaty wygasają bez
możliwości odnowienia, a~usługi statusu (CRL/OCSP) przestają być utrzymywane.
Weryfikacja długoterminowa wymaga wówczas dodatkowej infrastruktury
znakowania czasem i~archiwizacji materiału walidacyjnego (podpisy w~formatach
długoterminowych), która w~praktyce rzadko obejmuje dyplomy --- stąd ocena
częściowa, nie negatywna.

Zasadniczy wniosek z~tabeli jest dla blockchaina niewygodny, lecz uczciwy. PKI
\emph{dorównuje} rozwiązaniu blockchainowemu w~dwóch pierwszych, najczęściej
przywoływanych zaletach: pozwala na niezależną weryfikację offline bez kontaktu
z~wystawcą i~po znikomym koszcie jednostkowym. Podpis cyfrowy pod dyplomem daje to
samo bez konieczności utrzymywania blockchaina. Argument „blockchain eliminuje
telefon do dziekanatu'' nie jest więc argumentem za \emph{blockchainem} --- jest
argumentem za \emph{kryptografią asymetryczną}, którą PKI dostarcza taniej.

\section{Czy blockchain tu wygrywa?}
\label{sec:gov-wniosek}

Różnica pojawia się dopiero w~dwóch ostatnich wierszach tabeli~\ref{tab:gov-value},
dotyczących nie pojedynczego dokumentu, lecz \emph{warstwy uprawnień do wydawania}.
W~modelu PKI o~tym, które certyfikaty wystawców są ważne, decyduje urząd certyfikacji.
Może on unieważnić lub wystawić certyfikat prywatnie, a~listy unieważnień (CRL/OCSP)
odzwierciedlają \emph{bieżący} stan --- nie dają publicznego, niemodyfikowalnego
zapisu tego, kto i~kiedy był uprawniony, ani nie chronią przed cichą wsteczną
zmianą tego zbioru. Weryfikator ufa, że urząd nie przepisał historii.

\paragraph{Certificate Transparency zamyka tę lukę --- prawie.} Powyższy argument
byłby rozstrzygający, gdyby publiczny log uprawnień dało się uzyskać wyłącznie
na blockchainie. Tak nie jest. Certificate Transparency~\cite{rfc6962} realizuje
dokładnie tę własność --- publiczny, wyłącznie dopisywalny log wystawionych
certyfikatów z~dowodami spójności --- na drzewach Merkle, bez łańcucha bloków,
produkcyjnie od ponad dekady i~po koszcie nieporównanie niższym niż utrzymanie
kontraktu. Czwarta i~piąta własność w~tabeli~\ref{tab:gov-value} przestają zatem
odróżniać blockchain od dojrzałej alternatywy. Uczciwe postawienie sprawy wymaga
przyznania, że \emph{większość} przewagi, którą praca zamierzała wykazać, jest
osiągalna w~modelu PKI+CT.

Pozostaje różnica wąska i~jedna, ujęta w~ostatnim wierszu tabeli. Log CT jest
prowadzony przez wyznaczony zbiór operatorów, a~o~tym, którym z~nich klient ufa,
rozstrzyga polityka dostawcy przeglądarki. Log może wobec różnych obserwatorów
przedstawić \emph{rozdwojony widok} (\emph{split view}): dwa niesprzeczne
wewnętrznie, lecz różne stany, z~których każdy ma poprawne dowody spójności.
Wykrycie takiego rozdwojenia wymaga, by obserwatorzy porównywali swoje widoki
między sobą. Mechanizm do tego przeznaczony był specyfikowany w~IETF jako
osobny dokument~\cite{trans-gossip}, lecz nigdy nie osiągnął statusu standardu
--- prace nad nim wygasły. W~praktyce spójność widoku opiera się na wymogu
zebrania poświadczeń od kilku niezależnych operatorów oraz na ich wzajemnym
audycie, a~więc na założeniu, że operatorzy nie działają w~zmowie.

Publiczny łańcuch bloków rozwiązuje ten sam problem inaczej: jednolitość widoku
jest u~niego własnością warstwy konsensusu, nie dodatkową warstwą protokołu.
Nie ma logu, którego stan mógłby się rozdwoić wobec dwóch weryfikatorów, bo
o~stanie rozstrzyga ten sam mechanizm, który utrzymuje sam łańcuch.

Marginalna wartość blockchaina w~tym zastosowaniu sprowadza się zatem do
własności węższej, niż wynikałoby z~porównania z~samym PKI: \emph{publicznego,
dopisywalnego, odpornego na manipulację logu tego, kto może wydawać dyplomy,
o~jednym widoku dla wszystkich weryfikatorów, bez zaufania do zbioru operatorów
logu}. Logu, którego nie da się
po cichu wstecznie przepisać. Model zarządzania z~sekcji~\ref{sec:gov-model} jest
tym, co tę własność urzeczywistnia --- bez timelocka i~progu sam blockchain tej
gwarancji by nie dawał, bo pojedynczy właściciel mógłby log przepisać w~granicach
swoich uprawnień.

\paragraph{Studium przypadku: uczelnia skali politechniki.} Skalę kosztu dobrze
obrazuje rachunek dla dużej uczelni technicznej wydającej rzędu 10~tys.\ dyplomów
rocznie. Publikacja partii na~Base kosztuje około $9\cdot10^{-4}$~USD niezależnie
od jej rozmiaru (tab.~\ref{tab:issue-tx}). Nawet przy kilkunastu partiach rocznie
--- na przykład osobnej partii na wydział i~semestr --- łączny koszt publikacji nie
przekracza dwóch centów. Po doliczeniu jednostkowych unieważnień
($1{,}1\cdot10^{-3}$~USD, tab.~\ref{tab:revoke}) oraz sporadycznych,
amortyzowanych zmian uprawnień (sekcja~\ref{sec:gov-koszt}) całkowity roczny koszt
on-chain pozostaje poniżej dziesięciu centów.

Koszt transakcyjny nie jest zatem czynnikiem decyzyjnym w~żadnym realistycznym
scenariuszu wdrożenia. Realne koszty leżą w~integracji z~systemami dziekanatowymi
i~w~utrzymaniu operacyjnym (rozdz.~\ref{ch:podsumowanie}), a~wybór między
rozważanymi podejściami rozstrzyga się na osi wymagań, nie kosztów.

Pozostaje pytanie o~proporcję. Wyodrębniona wyżej własność kosztuje --- zależnie
od sieci --- ułamki mikrocenta na dyplom (rozdz.~\ref{ch:l2}) plus jednorazowy
narzut 740\,008 gazu na zmianę uprawnień. Jest to koszt niski w~wartościach
bezwzględnych, lecz niezerowy --- i~ponoszony wobec dwóch alternatyw, z~których
każda jest dojrzalsza.

Praca nie rozstrzyga tej proporcji arbitralnie. Wniosek ma postać stopniowanej
rekomendacji, a~nie opowiedzenia się po jednej stronie.

\begin{description}
  \item[Samo PKI wystarcza,] gdy wymaganiem jest wyłącznie niezależna od wystawcy,
        tania weryfikacja dokumentu. Argument ,,blockchain eliminuje telefon do
        dziekanatu'' jest w~całości zaspokojony na tym poziomie.
  \item[PKI z~Certificate Transparency wystarcza,] gdy dodatkowym wymaganiem jest
        publiczna, historyczna audytowalność uprawnień do wydawania --- i~jest to
        rozwiązanie właściwe dla zdecydowanej większości realnych wdrożeń,
        tańsze i~lepiej zbadane niż blockchain.
  \item[Blockchain wygrywa wtedy i~tylko wtedy,] gdy do powyższego dochodzi
        wymaganie \emph{jednolitości widoku bez zaufania do operatorów logu}:
        gdy weryfikatorzy nie mogą przyjąć, że wyznaczony zbiór operatorów nie
        działa w~zmowie i~nie przedstawia im rozdwojonego widoku, ani nie mogą
        polegać na warstwie uzgadniania widoków, która w~ekosystemie CT nie
        została ustandaryzowana. Scenariuszem modelowym jest weryfikacja
        ponadinstytucjonalna i~transgraniczna, w~której żaden pojedynczy podmiot
        --- ani urząd certyfikacji, ani operator logu --- nie jest akceptowany
        przez wszystkie strony.
\end{description}

Trzy powyższe progi, a~zwłaszcza wyodrębnienie trzeciego, stanowią główny wkład
analityczny niniejszej pracy. Jest on dla tezy blockchainowej mniej korzystny,
niż zakładano na wstępie: uwzględnienie Certificate Transparency przenosi
większość deklarowanej przewagi na rozwiązanie bez łańcucha bloków i~zawęża
uzasadnione zastosowanie blockchaina do przypadku brzegowego. Przypadek ten
pozostaje jednak realny, a~jego cena transakcyjna --- znikoma w~porównaniu
z~kosztem sporu o~to, komu wolno wystawiać dyplomy.
